% !PRINTED 473
\weberpage{473}

This leads us, then, to study more closely the permutations of $m$ figures
\[
0, 1, 2 \ldots m-1
\]
and their properties in general.

If the figure 0 passes into $a_0$, 1 into $a_1$, 2 into $a_2$, and so on, $m-1$ into
$a_{m-1}$, then

\weq{4}{a_0, a_1, a_2 \ldots a_{m-1}}

must be some other arrangement of the figures $0, 1, 2 \ldots m-1$. The
transition from the one arrangement to the other we denote by
a symbol, such as

\weq{5}{\pi_a = \begin{pmatrix} 0, & 1, & 2 & \ldots & m-1 \\ a_0, & a_1, & a_2 & \ldots & a_{m-1} \end{pmatrix}}

and we call this transition a \gesperrt{permutation}.

The number of all possible permutations of $m$ elements is
\[
\Pi(m) = 1 \,.\, 2 \,.\, 3 \ldots m,
\]
among which the \gesperrt{identical permutation}
\[
\begin{pmatrix} 0, & 1, & 2 & \ldots & m-1 \\ 0, & 1, & 2 & \ldots & m-1 \end{pmatrix},
\]
which leaves every element in its place, is also counted.

The symbol (5) does not change its meaning if one arranges the individual
pairs
\[
\begin{array}{ccccc}
0, & 1, & 2, & \ldots & m-1 \\
a_0, & a_1, & a_2, & \ldots & a_{m-1}
\end{array}
\]
differently. If then
\[
b_0, b_1, b_2 \ldots b_{m-1}
\]
is any arrangement of the figures $0, 1, 2 \ldots m-1$, we may
also put, in place of (5):

\weq{6}{\pi_a = \begin{pmatrix} b_0, & b_1, & b_2 & \ldots & b_{m-1} \\ a_{b_0}, & a_{b_1}, & a_{b_2} & \ldots & a_{b_{m-1}} \end{pmatrix},}

which is entirely equivalent to (5).

From this the very important concept of the \gesperrt{composition of
permutations} can be derived. If we replace an index $a$ by an index
$b$, and then $b$ by $c$, the effect is the same as if $a$ were replaced
at once by $c$. If, then, we carry out two permutations one after
the other, the result is a definite permutation,
which we can represent in the following way. If
%CONT>>
